\documentclass[11pt,a4paper]{article}
\usepackage[utf8]{inputenc}
\usepackage[T1]{fontenc}
\usepackage[brazil]{babel}
\usepackage{graphicx}
\usepackage{geometry}
\geometry{margin=2.5cm}

\title{\textbf{Aula 19 --- CNN no dataset cats\_vs\_dogs}}
\author{Breno Porto Pinheiro do Prado \quad RA 185196 \\ Guilherme Bento Ramos \quad RA 185226}
\date{}

\begin{document}
\maketitle

O exerc\'icio consiste em aplicar uma Rede Neural Convolucional (CNN) a um dataset de imagens do TensorFlow. Foi escolhido o dataset \texttt{cats\_vs\_dogs}, cujo objetivo \'e classificar imagens entre gatos e cachorros.

\medskip

O conjunto de dados foi dividido em 80\% para treino e 20\% para valida\c{c}\~ao. As imagens foram redimensionadas para $150\times150$ pixels e os valores dos pixels normalizados para o intervalo $[0,\,1]$. A rede foi constru\'ida com tr\^es blocos de camada convolucional seguida de pooling, uma camada densa de 64 neur\^onios e uma sa\'ida com ativa\c{c}\~ao sigmoid, adequada para classifica\c{c}\~ao bin\'aria.

\section*{Resultado}

\begin{figure}[h]
\centering
\includegraphics[width=0.95\textwidth]{Pasted image.png}
\caption{Sa\'ida do terminal ao rodar o c\'odigo.}
\end{figure}

O modelo terminou o treino de 10 \'epocas com:

\begin{itemize}
    \item \textbf{Loss de valida\c{c}\~ao:} 0,6932
    \item \textbf{Acur\'acia de valida\c{c}\~ao:} 49,01\%
\end{itemize}

A acur\'acia ficou perto de 50\%, o que significa que o modelo praticamente chutou. Isso acontece porque uma CNN simples sem aumento de dados n\~ao consegue aprender padr\~oes suficientes para generalizar bem nesse dataset. Para melhorar, seria necess\'ario usar t\'ecnicas como \textit{data augmentation} ou uma rede pr\'e-treinada.

\bigskip
\noindent\fbox{\parbox{\dimexpr\textwidth-2\fboxsep-2\fboxrule\relax}{\small
\textbf{Declara\c{c}\~ao de uso de IA:} este trabalho contou com aux\'ilio do assistente Claude Sonnet 4.6 (Anthropic) para redigir o c\'odigo e o relat\'orio.}}

\end{document}
